\subsection{Безразмерная форма уравнений термоупругости}

Для выявления характерных масштабов и определения доминирующих физических механизмов введём безразмерные переменные. В качестве характернXых масштабов выберем:
характерный линейный размер области $L$, характерную скорость распространения упругих волн $c$, характерную амплитуду смещений $U$ и характерное изменение температуры $\Delta T$.

Введём безразмерные величины:
\begin{equation}
x = L \tilde{x}, \quad
t = \frac{L}{c} \tilde{t}, \quad
u = U \tilde{u}, \quad
T = \Delta T \tilde{T}.
\end{equation}

Подставляя данные выражения в систему уравнений термоупругости и переходя к безразмерным переменным, получаем:

\begin{equation}
\frac{\partial^2 \tilde{u}_i}{\partial \tilde{t}^2}
=
\frac{\partial}{\partial \tilde{x}_j}
\left[
\tilde{\lambda}
\delta_{ij}
\frac{\partial \tilde{u}_k}{\partial \tilde{x}_k}
+
\tilde{\mu}
\left(
\frac{\partial \tilde{u}_i}{\partial \tilde{x}_j}
+
\frac{\partial \tilde{u}_j}{\partial \tilde{x}_i}
\right)
\right]
-
\beta \frac{\partial \tilde{T}}{\partial \tilde{x}_i},
\end{equation}

\begin{equation}
\frac{\partial \tilde{T}}{\partial \tilde{t}}
-
\frac{1}{\mathrm{Fo}}
\nabla^2 \tilde{T}
=
\eta
\frac{\partial}{\partial \tilde{t}}
\left(
\frac{\partial \tilde{u}_k}{\partial \tilde{x}_k}
\right).
\end{equation}

Здесь введены следующие безразмерные параметры:

\begin{itemize}
  \item
  \[
  \beta = \frac{\gamma \Delta T}{\rho c^2 U},
  \]
  — безразмерный параметр термоупругой связи, характеризующий влияние температурных возмущений на упругую динамику;

  \item
  \[
  \mathrm{Fo} = \frac{k}{C c L},
  \]
  — аналог числа Фурье, определяющий соотношение характерных времён теплопереноса и распространения упругих волн;

  \item
  \[
  \eta = \frac{\gamma T_0 U}{C \Delta T},
  \]
  — параметр, характеризующий интенсивность термомеханического источника в уравнении теплопереноса.
\end{itemize}

Анализ полученных безразмерных параметров позволяет выделить различные режимы распространения термоупругих волн. В частности, при малых значениях параметра $\mathrm{Fo}$ тепловые эффекты носят квазистатический характер, тогда как при $\mathrm{Fo} \sim 1$ наблюдается сильная связность упругих и тепловых процессов. Значения параметра $\beta$ определяют степень влияния температурных неоднородностей на динамику упругих волн.
